\documentclass[tikz,border=6pt]{standalone}
% 공통 프리앰블 — PM Day1 교재 그림 (TikZ standalone)
\usepackage{fontspec}
\setmainfont{Apple SD Gothic Neo}
\setsansfont{Apple SD Gothic Neo}
\setmonofont{Menlo}
\usepackage{amssymb}
\usepackage{tikz}
\usetikzlibrary{arrows.meta,positioning,calc,shapes.geometric,fit,backgrounds,matrix,decorations.pathreplacing}
\definecolor{navy}{HTML}{1F3A5F}
\definecolor{teal}{HTML}{2A7F62}
\definecolor{rust}{HTML}{C8553D}
\definecolor{sand}{HTML}{E8E1D5}
\definecolor{ink}{HTML}{222222}
\definecolor{grey}{HTML}{7A7A7A}
\definecolor{lightgrey}{HTML}{EFEFEF}
\definecolor{navylight}{HTML}{DCE6F2}
\definecolor{teallight}{HTML}{DDF0E8}
\definecolor{rustlight}{HTML}{F6E0DA}
\tikzset{
  every node/.style={font=\small, text=ink},
  box/.style={draw=navy, thick, rounded corners=2pt, align=center, inner sep=5pt, fill=white},
  boxfill/.style={box, fill=navylight},
  boxteal/.style={draw=teal, thick, rounded corners=2pt, align=center, inner sep=5pt, fill=teallight},
  boxrust/.style={draw=rust, thick, rounded corners=2pt, align=center, inner sep=5pt, fill=rustlight},
  boxgrey/.style={draw=grey, thick, rounded corners=2pt, align=center, inner sep=5pt, fill=lightgrey},
  arr/.style={-{Stealth[length=2.5mm]}, thick, navy},
  arrteal/.style={-{Stealth[length=2.5mm]}, thick, teal},
  arrrust/.style={-{Stealth[length=2.5mm]}, thick, rust},
  arrgrey/.style={-{Stealth[length=2.5mm]}, thick, grey},
  lbl/.style={font=\footnotesize, text=grey, align=center},
  src/.style={font=\scriptsize, text=grey, align=left},
  title/.style={font=\normalsize\bfseries, text=navy, align=center},
}

\begin{document}
\begin{tikzpicture}
\def\S{52mm}
% 네 도메인 (Snowden 배치: 우측 = 질서, 좌측 = 무질서)
\node[boxteal, minimum width=\S, minimum height=\S, anchor=south east, text width=46mm, align=left, font=\footnotesize] (cx) at (-2mm,2mm) {\textbf{Complex (복잡)}\\인과관계는 사후에만 알 수 있다\\[3pt]\textbf{probe → sense → respond}\\안전한 실험(safe-to-fail probes), 창발\\[4pt]{\scriptsize P1 예: "가맹점주가 실시간 판매 데이터 연동에 어떻게 반응할 것인가" → 파일럿 12개점이라는 probe를 먼저 던진다}};
\node[boxfill, minimum width=\S, minimum height=\S, anchor=south west, text width=46mm, align=left, font=\footnotesize] (cd) at (2mm,2mm) {\textbf{Complicated (난해)}\\인과관계는 존재하지만 분석이 필요하다\\[3pt]\textbf{sense → analyze → respond}\\전문가 분석, 좋은 사례(good practice)\\[4pt]{\scriptsize P1 예: "SAP 애드온으로 ECC와 연결한다" → 전문가가 분석하면 답이 나온다}};
\node[boxrust, minimum width=\S, minimum height=\S, anchor=north east, text width=46mm, align=left, font=\footnotesize] (ch) at (-2mm,-2mm) {\textbf{Chaotic (혼돈)}\\인과관계가 없다\\[3pt]\textbf{act → sense → respond}\\일단 행동해 안정화, 분석은 나중에\\[4pt]{\scriptsize P1 예: "OD-POS 서버 OS 지원 종료 후 보안 사고가 나면" → 일단 행동해 안정화}};
\node[box, draw=grey, fill=lightgrey, minimum width=\S, minimum height=\S, anchor=north west, text width=46mm, align=left, font=\footnotesize] (cl) at (2mm,-2mm) {\textbf{Clear / Obvious (명백)}\\인과관계가 명백하고 누구나 안다\\[3pt]\textbf{sense → categorize → respond}\\모범 사례(best practice) 적용\\[4pt]{\scriptsize 예: 표준 절차가 그대로 적용되는 반복 작업}};
% 중앙 disorder
\node[circle, draw=grey, fill=white, thick, minimum size=17mm, font=\scriptsize, align=center] at (0,0) {Disorder\\(무질서)\\어느 도메인인지\\모르는 상태};
% 상단 라벨
\node[lbl, anchor=south, font=\small\bfseries, text=navy] at (-\S/2-2mm,\S+3mm) {비질서(unordered) — 창발적 실무};
\node[lbl, anchor=south, font=\small\bfseries, text=navy] at (\S/2+2mm,\S+3mm) {질서(ordered) — 인과관계 분석 가능};
% 하단 경고
\node[box, draw=rust, fill=rustlight, text width=104mm, font=\footnotesize, anchor=north] at (0,-\S-6mm) {\textbf{오용 경고} — "우리 프로젝트는 complex하니까 애자일"은 틀렸다.\\① Cynefin의 처방은 방법론 선택이 아니라 \textbf{의사결정 순서의 전환}이다.\\② 도메인 판정은 \textbf{사후적}이다 — 사전에 프로젝트 전체에 라벨을 붙이지 말고, \textbf{결정 하나하나}에 적용한다.\\하나의 프로젝트 안에 네 도메인이 공존한다.};
\node[src, anchor=north west] at (-\S-2mm,-\S-30mm) {출처: Snowden \& Boone (2007), A Leader's Framework for Decision Making, Harvard Business Review를 바탕으로 재구성. P1 예시는 교재 2.6절.};
\end{tikzpicture}
\end{document}
